\documentclass{article}
\usepackage[left=2cm, right=2cm, top=2cm, bottom=2cm]{geometry}
\usepackage{graphicx}
\usepackage{amsmath}
\usepackage{amssymb}
\usepackage{amsthm}
\usepackage{fancyhdr}
\usepackage{verbatim}
\usepackage{listings}
\usepackage{xcolor}
\usepackage{pdfpages}
\usepackage{cancel}
\usepackage{physics}
\usepackage{esint}
\usepackage{braket}

\lstset{
    language=C++,
    basicstyle=\ttfamily\footnotesize,
    keywordstyle=\color{blue},
    commentstyle=\color{green},
    stringstyle=\color{red},
    numbers=left,
    numberstyle=\tiny\color{gray},
    frame=single,
    breaklines=true,
    captionpos=b,
}


\title{MTH 446: Lecture \# 7}
\author{Cliff Sun}

\newtheorem{theorem}{Theorem}[section]
\newtheorem{lemma}[theorem]{Lemma}
\newtheorem{definition}[theorem]{Definition}
\newtheorem{conjecture}[theorem]{Conjecture}
\newtheorem{proposition}[theorem]{Proposition}
\newtheorem{corollary}[theorem]{Corollary}
\newtheorem{one minute paper}[theorem]{One Minute Paper}

\pagestyle{fancy}
\lhead{\textbf{\thepage}\ \ \nouppercase{\rightmark}}
\chead{MTH 446: Lecture \# 7}
\rhead{Cliff Sun}

\begin{document}

\maketitle

\section*{Lecture Span}
\begin{enumerate}
    \item Topology of the complex plane
\end{enumerate}

\begin{definition}
    A \underline{polygonal} line is formed by a finite number of line intervals, joined end to end 
\end{definition}

\begin{definition}
    A non-empty set $A$ is called \underline{(polygonally) connected} if every two points of $A$ can be connected by a polygonal line that lies completely in $A$. 
    In otherwords, you can draw a continuous line from $z_0$ to $z_1$ for every $z_0,z_1 \in A$. 
\end{definition}

\begin{definition}
    A non-empty open set is called \underline{domain} if it is connected. 
\end{definition}

\begin{definition}
    A domain together with some, none, or all of its boundary points is called a \underline{region}. 
\end{definition}

\begin{definition}
    Let $A \subseteq \mathbb{C}$ and $z_0 \in \mathbb{C}$. The point $z_0$ is called an \underline{accumulation point} of $A$ if 
    \begin{equation}
        \forall \epsilon > 0, \;\; \overset{\circ}{D}_{\epsilon}(z_0) \cap A \neq \emptyset
    \end{equation}
    Note, $\overset{\circ}{D}$ is the \textit{deleted $\epsilon$ neighorborhood!!!} Therefore, every $z_0$ is an accumlation point if every deleted $\epsilon$ neighorborhood of $z_0$ conntains at least one point of $A$. 
\end{definition}

A single point is not an accumulation point!

\subsection*{Example}

Let $A = \{i\}$. Then prove that $i \in \partial A$. and $i$ is not an accumulation point. 

\begin{enumerate}
    \item We first prove that $i \in \partial A$. In particular, $D_{\epsilon}(i) \cap A \neq \emptyset$ and $D_{\epsilon}(i) \cap \mathbb{C}/A \neq \emptyset$. Therefore, $i \in \partial A$.
    \item $i$ is not an accumulation point, because $\overset{\circ}{D}_{\epsilon}(i) \cap A = \emptyset$. 
\end{enumerate}

\subsection*{Example}

Suppose 

\begin{equation}
    A = \{z \in \mathbb{C}/\{0\}: 0 \leq \arg(z)  \pi/2\}
\end{equation}

Then, prove that $z=0$ is an accumulation point. 

\begin{proof}
    Draw a deleted $\epsilon$ neighorborhood at $0$. From the sketch, we see that $\overset{\circ}{D}_{\epsilon}(0) \cap A \neq \emptyset$. Therefore, $0$ is an accumulation point of $A$. 
\end{proof}

\subsection*{Example}

Let

\begin{equation}
    A = \{z_n = i^n: n = 1,2,3,\dots\}
\end{equation}

Find the set of all accumulation points.

\begin{proof}
    It is obvious that $A = \{1,-1,i,-i\}$. Then let $\epsilon_1 = \min_i\{|z_i - i|\}$, then for $\epsilon < \epsilon_1$, we have that $\overset{\circ}{D}_{\epsilon}(z_0) \cap A = \emptyset$. Therefore, $i$ is not an accumulation point. By symmetry, it follows that all other points of $A$ are not accumulation points.  
    Next, let $z_0 \notin A$. Then, let $\epsilon_2 = \min_i\{|z_i - z_0|\}$, then choose $\epsilon < \epsilon_2$, then $D_{\epsilon}(z_0) \cap A = \emptyset$. Therefore, $z_0$ is not an accumulation point. 
\end{proof}

Therefore, the set of all accumulation points of is $\emptyset$.

\begin{definition}
    A point $z_0 \in F$ is an \underline{isolated point} of $F$ if $$\exists \epsilon > 0: \overset{\circ}{D}_{\epsilon}(z_0) \cap F = \emptyset$$
    For example, all points of the set $A= \{i, 1, -1, -i\}$. 
\end{definition}

It follows that every point of $F$ is either an accumulation point of $F$ or it is an isolated point of $F$. 

\section*{Functions of complex variables}

Let $D \subseteq \mathbb{C}$ be non empty. Then $w = f(z)$, where $z \in D$ and $w \in \mathbb{C}$ is called the \underline{function of the complex variable}, i.e. 

\begin{equation}
    f: D \subseteq \mathbb{C} \rightarrow \mathbb{C}
\end{equation}

$D$ is called the domain of definition. 

\subsection*{Example}

\begin{equation}
    f(z) - \frac{1}{1 + z^2}
\end{equation}

Then the domain of $f$ is $\mathbb{C}/\{i, -i\}$.

\subsection*{Example}

$f(z) = |z|^2$. Then domain of $f$ is $\mathbb{C}$. 

\subsection*{Example}

Consider 

\begin{equation}
    P_n(z) = \sum_i a_i x^i
\end{equation}

With $a_i \in \mathbb{C}$ for all $i$. Then $P_n(z)$ is called a \underline{polynomial of degree $n$}. Then, the domain of $P_n$ is $\mathbb{C}$. 

\subsection*{Example}

\begin{equation}
    f(z) = \arg(1/z)
\end{equation}

Then, the domain of $f$ is $\mathbb{C}/\{0\}$. 

Next, set $u(x,y) = \Re(f(z)) = \Re(f(x + iy))$ and $S(x,y) = \Im(f(z))$. Then 

\begin{equation}
    f(z) = u(x,y) + iS(x,y)
\end{equation}

\subsection*{Example}

Let $f(z) = z^2$. Then $f(z) = (x + iy)^2 = zx^2 + 2xyi - y^2 = (x^2 - y^2) + 2xyi$. Therefore, $u = x^2 - y^2$ and $s = 2xy$. 

\subsection*{Example}

Let $f(z) = x^2 + 2y + 5xi$. Write $f$ in terms of $z$. Let $x = (z + \overline{z})/2$ and $y = (z - \overline{z})/2i$. Then plug these substitutions into the function. 

Next time, we will discuss the polar coordinates. 

\end{document}